\section{Conclusion}
\label{sec:conclusion}

\subsection{Summary of Contributions}
\label{subsec:concl_summary}

This thesis presented a GPU-accelerated computational investigation into the role of spatial heterogeneity, mass-conserving fluid mechanics, and curiosity-driven exploration within the continuous cellular automata framework of Flow Lenia. By re-engineering the continuous physics engine into native functional JAX, we eliminated artificial coordinate damping and evaluated open-ended evolutionary dynamics across structured environments.

The primary contributions of this work are summarized as follows:

\begin{enumerate}
    \item \textbf{Conservative Physics Engine in Functional JAX:} By formulating continuous convolutions via two-dimensional spectral Fourier transforms and implementing semi-Lagrangian bilinear reintegration tracking \parencite{moroz2020reintegration, plantec2025flowlenia}, our framework preserves global mass within floating-point round-off limits. This eliminates artificial coordinate drag and provides an elastomeric, soft-matter fluid substrate executing at over 33,000 steps per second.
    \item \textbf{Curiosity-Driven Automated Discovery and Quality Gating:} We developed an Intrinsically Motivated Goal Exploration Process (IMGEP) operating across a standardized three-dimensional metric space ($[\Motility, \EvoActivity, \Complexity]$). Evaluated through a sequential five-stage morphological filter, IMGEP achieved a 44.0\% gate pass rate in our benchmark ($n=3$ seeds) and discovered diverse families of swimming solitons while mapping critical failure modes such as stochastic seed sensitivity and the solidity-motility trade-off.
    \item \textbf{Biomechanical Characterization under Spatial Heterogeneity:} 
    \begin{itemize}
        \item In directional chemotaxis assays, we mapped a cohesion-fission phase transition in the explored parameter space, distinguishing unitary droplet locomotion ($96.5\%$ core preservation) from shear-induced amoeboid division.
        \item In a two-chamber aperture constriction assay, we measured the aperture transmission coefficient $T(W)$, showing that soft-bodied solitons exhibit a sigmoidal transmission profile with an empirical necking transition interval near $W \approx 28\text{--}32\text{ pixels}$.
    \end{itemize}
    \item \textbf{Dynamic Niche Construction and Ecological Dynamics:} Coupling local growth potentials to dynamic resource depletion eliminated static crystal attractors in our evaluated sample ($0.0\%$ stationary states), compelling continuous cyclic foraging trajectories. Across 22,500 continuous simulation steps in a multi-chamber arena benchmark, categorical Gumbel-Max flux sampling prevented parameter homogenization and preserved distinct territorial lineages.
\end{enumerate}

\subsection{Answers to Research Questions}
\label{subsec:concl_answers}

Returning to the central research question:

\begin{quote}
\textbf{Central Research Question:} \textit{To what extent does structured spatial heterogeneity drive phenotypic adaptation, structural plasticity, and behavioral divergence in continuous mass-conserving cellular automata?}
\end{quote}

Our empirical findings indicate that structured spatial heterogeneity acts as an essential catalyst for complex dynamics in continuous cellular automata:
\begin{itemize}
    \item In homogeneous, isotropic spaces, continuous systems suffer from evolutionary stagnation: gliders encounter zero environmental resistance, and populations converge to static or trivial periodic attractors.
    \item Introducing spatial friction (rigid geometric constrictions, depletable resource fields, and external chemotactic gradients) destabilizes static equilibria in the tested regimes, tests droplet elasticity, forces adaptive deformation, and sustains long-term migratory and territorial dynamics.
\end{itemize}

Addressing the specific sub-questions:
\begin{itemize}
    \item \textbf{Sub-Question 1 (Curiosity Discovery):} IMGEP effectively circumvents the parameter desert in our benchmark by expanding the behavioral frontier in continuous metric space. Multi-seed ensemble evaluation mitigates seed sensitivity, while non-objective goal babbling prevents premature convergence to the low-complexity attractors that trap scalar objective optimization.
    \item \textbf{Sub-Question 2 (Aperture Transmission):} Soliton transmission through narrow slits is governed by a balance between external gradient pull and internal cohesive surface tension, establishing an aperture interval below which passage is mechanically obstructed.
    \item \textbf{Sub-Question 3 (Niche Construction and Scalability):} Dynamic resource depletion introduces a self-generated negative feedback loop that prevents organisms from resting, driving continuous trail following and multi-species ecological stability across large spatial canvases.
\end{itemize}

\subsection{Epistemological Synthesis and Computational Limits}
\label{subsec:concl_epistemology}

This research established clear epistemological boundaries regarding Artificial Life in continuous cellular fields:
\begin{itemize}
    \item \textbf{Physical Droplets vs. Cognitive Agents:} Solitons in Flow Lenia are continuous dissipative wave structures governed by partial differential equations. They exhibit soft-bodied elastomeric physics rather than symbolic information processing or cognitive decision-making.
    \item \textbf{Demarcation of Spatial Routing:} When continuous solitons traverse complex mazes or solve routing problems, global path planning is governed by external geodesic potential fields, while Flow Lenia supplies the continuous, deformable physical medium.
    \item \textbf{Missing Physical Constraints:} The absence of impenetrable physical membranes and thermodynamic energy accounting limits the emergence of true open-ended evolutionary hierarchies. In a medium where mass updates carry zero energetic cost, selective pressure for metabolic efficiency is fundamentally absent.
\end{itemize}

\subsection{Future Outlook: Towards Hybrid Architectures}
\label{subsec:concl_future_work}

To bridge the remaining gap in artificial open-ended evolution, future research should focus on three promising avenues:
\begin{enumerate}
    \item \textbf{Thermodynamic Energy Accounting:} Formulating explicit local energy budgets where mass advection, state updates, and kernel convolutions consume finite metabolic resources, penalizing stationary dissipation and rewarding efficient foraging.
    \item \textbf{Physical Enclosed Membranes:} Developing phase-field or level-set models that support impermeable, deformable boundaries capable of maintaining chemical isolation against external mixing.
    \item \textbf{Hybrid HPC Architectures:} Combining GPU-accelerated continuous field equations for underlying fluid mechanics with asynchronous, discrete agent-based processing for downward causal decision-making, resolving the parallelization paradox.
\end{enumerate}

In conclusion, while continuous cellular automata do not spontaneously generate biological cognition, coupling mass-conserving fluid mechanics with structured spatial heterogeneity provides a mathematically sound, highly scalable foundation for modeling soft-matter physics, emergent biomechanics, and decentralized artificial life.